\documentclass[reprint, amsmath, amssymb, aps, prd]{revtex4-2}

% --- PACKAGES ---
\usepackage[utf8]{inputenc}
\usepackage[T1]{fontenc}
\usepackage{graphicx} 
\usepackage{hyperref} 
\usepackage{bm}       
\usepackage{booktabs} 

\begin{document}

\title{A Phase-Ordered Pre-Geometric Projection Framework}
\author{Richard Fuoco}
\affiliation{Independent Researcher, Saint-Basile-le-Grand, QC, Canada}
\date{February 21, 2026}

\begin{abstract}
We propose a relational/algebraic substrate and a candidate finite-capacity projection program. Current exact toy models recover interaction locality encoded in selected chain/grid Hamiltonians, choose one- and two-dimensional finite embeddings, and test graph clock constraints. Raw relative entropy fails as a standalone linear source; a microscopic KMS-energy density remains a finite-model candidate. General Relativity, Lorentz recovery, a three-dimensional continuum limit, and singularity resolution are matching goals or conjectures, not derived results.
\end{abstract}

\maketitle

\paragraph{Scientific-status note (August 9, 2026).}
This is a theoretical submission draft. Its unique projection,
$3+1$-dimensional recovery, GR/QFT closure, Lorentz recovery, and singularity
claims are hypotheses or requirements, not completed results. The companion
code currently supports finite exact tests of partition ranking, blind recovery
of encoded chain/grid locality, MDS dimension selection, a numerical graph
clock constraint, local SPD embedding fits, and a 2D angle-deficit proxy. It
does not implement QCMI filtering, intrinsic Regge/Einstein closure, or a
validated physical Araki/KMS source. Perturbative tests reject
raw relative entropy as a standalone linear source.

% ----------------------------------------------------------------
\section{Reviewer-facing summary}
% ----------------------------------------------------------------
\subsection{What this framework commits to}
\begin{itemize}
    \item \textbf{Substrate:} an atemporal, pre-geometric relational/algebraic structure $(\mathcal{A}, \omega)$ with a compact internal symmetry and a preferred \textbf{phase/action ordering flow} $\sigma_s$ (order parameter $s$, not time).
    \item \textbf{Projection:} a \textbf{unique and necessary} physical map $\Pi$ from substrate to effective descriptions. $\Pi$ is defined operationally by (i) a stability-selected coarse-graining into finite-capacity ``cells'', (ii) correlation-defined locality, (iii) an embedding procedure into 3D, and (iv) a phase-order-to-clock-time mapping.
    \item \textbf{Finite distinguishability:} any finite projected region has bounded distinguishable information, scaling with \textbf{boundary area} (Planck-area-like scale $\ell_*^2$), interpreted as a \textbf{relational encoding boundary} rather than a geometric boundary in the substrate.
    \item \textbf{Emergence:} discrete quantum ``types'' arise from stable representation content under projection constraints; effective geometry arises from correlation structure and the clock mapping, not from substrate curvature.
\end{itemize}

\subsection{What this framework does \textit{not} claim (yet)}
\begin{itemize}
    \item It does \textbf{not} assume a complete derivation of the Standard Model spectrum, coupling constants, or cosmological parameters in this draft.
    \item It does \textbf{not} assert superluminal signaling. ``Junction access'' is defined as a constrained open module with parameters required to match existing no-signaling constraints in all empirically accessed regimes.
    \item It does \textbf{not} claim a closed-form fundamental ``constitutive law'' $g = F(\rho, \ldots)$; the effective metric is defined by the explicit projection construction (graph distances $\rightarrow$ embedding $\rightarrow$ reconstructed $h_{ab}$ and $d\tau$ mapping).
\end{itemize}

\subsection{Scientific status and evaluation}
\begin{itemize}
    \item The framework is \textbf{empirically anchored} by matching requirements: in accessible regimes, the induced effective description must reproduce (to stated tolerance) standard GR tests and standard quantum statistics.
    \item The framework becomes \textbf{strictly falsifiable} when it asserts: (i) universal area-law capacity bounds outside their known domain, (ii) saturation/no-singularity behavior in regimes where GR predicts divergence, or (iii) any nonzero operational ``junction access'' parameter in regimes already constrained by experiment.
    \item Absent additional predictive commitments, the framework is evaluated by: (a) internal consistency (no hidden time/geometry in substrate), (b) uniqueness/minimality of $\Pi$ under stated selection principles, and (c) whether GR/QFT emerge as stable effective closures with a small, auditable parameter budget.
\end{itemize}

% ----------------------------------------------------------------
\section{Motivation and scope}
% ----------------------------------------------------------------
This document introduces a conceptual framework intended for comprehensive evaluation and iterative refinement. The goal is to provide a minimal set of primitives and postulates from which (i) an effective 3+1 spacetime description, (ii) quantum statistical behavior, and (iii) General Relativity in tested regimes can be recovered as projection-level physics. The framework is explicitly pre-geometric at the substrate level and treats time as emergent from a phase/action ordering structure.

Scope control: the present draft focuses on structural definitions, matching requirements, and evaluation criteria. It does not claim a completed derivation of the Standard Model or the Born rule; instead it identifies the minimal mathematical objects needed to attempt such derivations.

% ----------------------------------------------------------------
\section{Design constraints}
% ----------------------------------------------------------------
\begin{itemize}
    \item \textbf{No substrate time:} the substrate admits no fundamental temporal parameter.
    \item \textbf{No substrate spatial geometry:} the substrate is not a manifold with metric/curvature; geometric notions are projection outputs.
    \item \textbf{Analog intuition with bounded physics:} continuous symmetry/phase structure is allowed as an idealization, but no unbounded physical observables (no divergent densities; no infinite recursion depth).
    \item \textbf{Unique and necessary projection:} the mapping from substrate to projection is physical and not contingent among many alternatives.
    \item \textbf{Finite distinguishability:} any finite projected region has a bounded number of distinguishable states; capacity scales with boundary area (area law) via a relational encoding boundary.
    \item \textbf{Empirical recovery:} in currently tested regimes the framework must reproduce standard GR and quantum predictions to within experimental bounds.
    \item \textbf{FTL signaling remains an open module:} the existence of operational superluminal channels is not assumed; it is parameterized and constrained.
\end{itemize}

% ----------------------------------------------------------------
\section{Core postulates (axioms)}
% ----------------------------------------------------------------

\subsection{Substrate postulates}
\textbf{S1 (Atemporal substrate).} There exists a substrate $\mathcal{S}$ that is not spacetime and has no intrinsic time parameter.  

\textbf{S2 (Pre-geometric substrate).} The substrate carries relational/algebraic structure but no spatial metric, curvature, or geometric topology in the usual sense.  

\textbf{S3 (Compact internal symmetry).} The substrate admits a compact internal SU(2)-like symmetry acting on its relational degrees of freedom.  

\textbf{S4 (Phase/action ordering generator).} A distinguished generator induces an ordering structure (phase/action order) that is not time but is used by projection to construct time-order.

\subsection{Projection postulates}
\textbf{P1 (Equivalence-Class Uniqueness).} There exists a physical projection map $\Pi$ from substrate structure to projection-level effective physics. While the specific microscopic cell net may only be unique up to spontaneous symmetry breaking (attractor basin selection, analogous to a gauge choice), the \textit{macroscopic predictions and geometric invariants} output by $\Pi$ are uniquely determined, establishing equivalence-class uniqueness.

\textbf{P2 (Emergent spatial structure; dimension as output).} $\Pi$ yields an emergent metric space with a manifold-like realization of some effective dimension $D^*$; in the empirically accessed regime of our universe, $D^* \approx 3$ (treated as an output/matching condition rather than an assumed input).  

\textbf{P3 (Emergent time-order).} $\Pi$ yields an objective time-order in projection as a metric on phase/action order; non-conscious systems inherit this time-order.  

\textbf{P4 (Finite distinguishability; boundary-capacity scaling).} For any finite projected region $R$, physically distinguishable information is bounded by a boundary \textit{cut-capacity} functional (defined on the correlation graph). In manifold-like regimes this scaling is expected to reproduce an area-law.  

\textbf{P5 (Finite-capacity regularization conjecture).} Projection-level observables may saturate rather than diverge. This is not yet demonstrated dynamically and does not imply causal or geodesic completeness.

\subsection{Emergence postulates}
\textbf{E1 (Quantum discreteness from representations).} Discrete quantum types arise from stable representation content of the internal symmetry under projection constraints, not from fundamental substrate discreteness.  

\textbf{E2 (Geometry as projection output).} The projection is intended to construct an effective metric from correlation distances, an embedding/local-metric diagnostic, and a clock mapping. The current code implements finite prototypes of those pieces but not a complete spacetime metric or tensor dynamics.

\textbf{E4 (Lorentz recovery conjecture).} SU(2)-equivariance constrains internal/spatial rotations only; it does not implement boosts or establish Lorentz covariance. Infrared recovery remains conjectural until dispersion, common-speed, anisotropy, boost, and renormalization-group calculations are supplied.

\textbf{J (module) (Junction accessibility).} In all known regimes, operations available to agents obey no-signaling in local marginals. Additional junction access is treated as a constrained open module.

% ----------------------------------------------------------------
\section{Mathematical primitives}
% ----------------------------------------------------------------

\subsection{Substrate representation (analog-friendly)}
Represent the substrate as a pair:
\begin{equation}
\mathcal{S} = (\mathcal{A}, \omega)
\end{equation}
where $\mathcal{A}$ is an algebra of relational observables and $\omega$ is a state on $\mathcal{A}$. No manifold, metric, or time parameter is assumed at this level.

\textbf{Analog-friendly choice.} To support continuous internal symmetry/phase structure while preserving \textit{finite physical realizability} (via projection-level finite distinguishability), take $\mathcal{A}$ to be a separable operator algebra that is well-approximated by finite matrix algebras at any finite resolution. A convenient template is:
\begin{equation}
\mathcal{A} = \mathcal{A}_{\text{rel}} \otimes \mathcal{A}_F
\end{equation}
where $\mathcal{A}_{\text{rel}}$ is the ``relational bulk'' algebra (pre-geometric), chosen so that it can be approximated by an increasing sequence of finite-dimensional matrix algebras:
\begin{equation}
\mathcal{A}_{\text{rel}} \approx \text{closure}\left( \bigcup_n M_{k_n}(\mathbb{C}) \right)
\end{equation}
This supports the \textbf{analog} intuition (continuous symmetry/phase) while remaining compatible with the principle that \textit{only a finite number of states are physically distinguishable in any finite projected region} (Section 6).

$\mathcal{A}_F$ is an optional finite ``internal'' algebra used to encode gauge/representation structure in a purely algebraic way. A minimal candidate that naturally contains $U(1)$, $SU(2)$, and $SU(3)$-type unitary structure is:
\begin{equation}
\mathcal{A}_F = \mathbb{C} \oplus \mathbb{H} \oplus M_3(\mathbb{C})
\end{equation}
with $\mathbb{H}$ the quaternions.

\textbf{Remark (no infinities physically instantiated).} Continuous groups and infinite-dimensional algebras are treated as \textit{idealized descriptions}: the framework's finiteness claim is that no \textit{projection-level observable} diverges and that any finite region has finite operational capacity (area-law distinguishability), not that the descriptive mathematics cannot employ limits.

\subsection{Internal symmetry action}
Let $\alpha$ be an action of the internal symmetry group on $\mathcal{A}$ as *-automorphisms:
\begin{equation}
\alpha: \text{SU}(2) \to \text{Aut}(\mathcal{A})
\end{equation}
Symmetry invariants are quantities unchanged under $\alpha_U$ for all $U \in \text{SU}(2)$.

\subsection{Phase/action ordering generator}
The substrate has no time, but it provides an \textbf{ordering structure} that projection uses to construct time-order.

\textbf{(i) Generator form (working placeholder).} Let $H \in \mathcal{A}$ be a distinguished self-adjoint element ($H = H^\dagger$) defining a one-parameter family of inner automorphisms:
\begin{equation}
\mathcal{A} \mapsto e^{isH} \mathcal{A} e^{-isH}
\end{equation}
The parameter $s$ is a phase/action \textbf{order parameter}, not substrate time.

\textbf{(ii) Canonical-flow form (parameter-reducing option).} Impose that the pair $(\mathcal{A}, \omega)$ determines a canonical one-parameter automorphism flow $\sigma_s$ on $\mathcal{A}$. Projection-time is then constructed from monotone ``distance'' along this flow. In this view, $H$ is not chosen freely; it is the generator associated with the canonical flow in a suitable representation.

\subsection{Projection map (explicit construction template)}
In this framework, the projection map $\Pi$ is not treated as an arbitrary ``fit function.'' Instead it is defined as a \textit{constrained construction} that turns substrate relational structure into:
\begin{itemize}
    \item an emergent 3D locality structure,
    \item an effective spatial metric,
    \item an objective time-order (from phase/action order),
    \item an encoding-density field and finite distinguishability bounds,
    \item effective subsystem states used for quantum predictions.
\end{itemize}

Mathematically, it is useful to write $\Pi$ as a composition of four stages:
\begin{equation}
\Pi = \Pi_{\text{time}} \circ \Pi_{\text{geom}} \circ \Pi_{\text{loc}} \circ \Pi_{\text{res}}
\end{equation}

\subsubsection{Stage 0: represent the substrate state (GNS representation)}
Given $(\mathcal{A}, \omega)$, choose the GNS triple $(\pi_\omega, \mathcal{H}_\omega, |\Omega\rangle)$ such that:
\begin{equation}
\omega(a) = \langle\Omega | \pi_\omega(a) | \Omega\rangle \quad \text{for all } a \in \mathcal{A}.
\end{equation}

\subsubsection{\texorpdfstring{$\Pi_{\text{res}}$}{Pi\_res}: resolution-limited coarse-graining via split-property funnels}
A projection must provide a notion of subsystems/regions, but subsystems are not substrate primitives. Furthermore, strict local algebras in Quantum Field Theory are typically \textbf{Type III factors} (entanglement is ubiquitous), meaning they do not admit pure states or density matrices, and cannot be factorized into tensor products $\mathcal{A}_R \otimes \mathcal{A}_R^c$.

To resolve this \textbf{Type III incompatibility} without abandoning the continuum, we define emergent ``cells'' explicitly as \textbf{split-property funnels} (Type I approximants across algebraic buffer zones).

\textbf{The Funnel Definition.}
Instead of a single algebra $\mathcal{A}_i$, the projection selects a \textbf{funnel structure} for each emergent site $i$:
\begin{equation}
\mathcal{A}_{i, \text{inner}} \subset N_i \subset \mathcal{A}_{i, \text{buffer}} \subset \mathcal{A}
\end{equation}
where $N_i$ is a \textbf{Type I factor} (isomorphic to $\mathcal{B}(\mathcal{H})$ or a finite matrix algebra $M_d(\mathbb{C})$ in toy models) that acts as a local ``distinguishable semantic shell'' between the inner core and the buffer zone.

The effective local algebra is identified with the Type I factor $N_i$. This guarantees the existence of well-defined local density matrices, entropies, and particle-like excitations.

\textbf{Coarse-Graining.}
The coarse-graining maps $E_i$ are conditional expectations onto these Type I factors:
\begin{equation}
\{\mathcal{A}_i \subset \pi_\omega(\mathcal{A})'' \}_{i \in V}, \quad \text{with} \quad \mathcal{A}_i \cong M_{d_i}(\mathbb{C}),
\end{equation}
together with conditional expectations (coarse-graining maps):
\begin{equation}
E_i : \pi_\omega(\mathcal{A})'' \to \mathcal{A}_i
\end{equation}
that are completely positive, unital, and idempotent ($E_i \circ E_i = E_i$).

For any finite set of cells $R \subset V$, define the induced effective algebra and reduced state:
\begin{equation}
\mathcal{A}_R \approx \bigotimes_{i\in R} \mathcal{A}_i, \quad E_R := \bigotimes_{i\in R} E_i, \quad \omega_R := \omega \circ E_R.
\end{equation}

\textbf{Finite distinguishability constraint (non-geometric form).}
To avoid circularity, the capacity bound is imposed \textbf{before} any geometric embedding. Let $\mathrm{Cap}(\partial R)$ denote a \textbf{cut-capacity} functional on the cell net:
\begin{equation}
\mathrm{Cap}(\partial R) := \sum_{i\in R, j\notin R} \kappa(I_{ij}),
\end{equation}
where $I_{ij}$ is the mutual information between cells $i$ and $j$, and $\kappa$ is a fixed monotone map to dimensionless edge capacities.

\textbf{Finite distinguishability constraint (Araki-relative form).}
For Type III compatibility, we eschew the naive von Neumann entropy. The bound is imposed on the \textbf{Araki relative entropy} of the region's state $\omega|_R$ against the canonical \textbf{KMS vacuum baseline} $\omega^{\text{vac}}|_R$:
\begin{equation}
S_{\text{Araki}}( \omega|_R \| \omega^{\text{vac}}|_R ) \leq \eta \cdot \mathrm{Cap}(\partial R)
\end{equation}
A geometric area scaling $\mathrm{Cap}(\partial R) \propto A_{\text{geo}}(\partial R)/\ell_*^2$ is treated as an emergent behavior in the low-stress embedding regime (Section~\ref{subsec:pi-geom}), not as an input constraint.

\subsubsection*{Selection Principle \texorpdfstring{$E$}{E}: symmetry-commuting, phase-flow stable cell net}
The choice of the effective cell algebras $\{\mathcal{A}_i\}$ and coarse-graining maps $\{E_i\}$ is fixed by a \textit{stability principle} tied to the phase/action ordering flow $\sigma_s$. Intuitively: the correct decomposition is the one for which phase/action ordering does not continually ``slosh'' information across emergent cell boundaries.

Let $\alpha_g$ denote the internal SU(2)-like symmetry action on $\mathcal{A}$. Let $\sigma_s$ denote the phase/action ordering flow. Define $\mathfrak{E}_{\text{adm}}$ as the set of families $\{E_i\}$ (and their associated cell algebras $\{\mathcal{A}_i\}$) satisfying all of:

1) \textbf{Finite-capacity cells.} Each cell algebra is finite-dimensional $\mathcal{A}_i \cong M_{d_i}(\mathbb{C})$, with either a fixed common dimension $d_i = d$ for all cells, or an explicit local bound $d_i \leq d_{\max}$. 

2) \textbf{Symmetry commutation (internal isotropy constraint).} Each coarse-graining map commutes with the SU(2)-like symmetry:
\begin{equation}
E_i \circ \alpha_g = \alpha_g \circ E_i \quad \text{for all } g \in \text{SU}(2) \text{ and all } i.
\end{equation}

\textbf{Lorentz status.} This constraint supplies a spatial-rotation equivariance condition in the toy coarse graining. The proposed custodial/spin-connection mechanism is unimplemented, so no conclusion about boosts, Lorentz-violating operators, or RG suppression follows.

3) \textbf{Information retention (anti-triviality constraint).} Impose a global bound on relative-entropy loss:
\begin{equation}
D( \omega \| \omega \circ E ) \leq \varepsilon,
\end{equation}

\textbf{Primary stability functional (phase-flow leakage).} For each admissible $\{E_i\}$, define the phase-flow leakage:
\begin{equation}
\mathcal{L}_{\text{leak}}(E) := \int ds \, w(s) \sum_i \| E_i \circ \sigma_s - \sigma_s \circ E_i \|^2,
\end{equation}
where $w(s) \geq 0$ is a weight over a chosen phase-order interval. Exact co-motion corresponds to $\mathcal{L}_{\text{leak}} = 0$.

\textbf{Optional tie-breaker (local information drift).} Define $\omega_s := \omega \circ \sigma_s$, and let $\rho_i(s)$ be the density operator corresponding to the reduced state $\omega_s \circ E_i$. For a small step $\delta$, define:
\begin{equation}
\mathcal{L}_{\text{drift}}(E) := \int ds \, w(s) \sum_i \frac{1}{\delta^2} D( \rho_i(s+\delta) \| \rho_i(s) ).
\end{equation}

\textbf{Cell selection as causal gradient flow.}
The cell net $E^* = \{E_i^*\}$ is \textbf{not} defined by a global argmin. Instead, it is defined as the \textbf{stable fixed-point (attractor)} of a local, causally-bounded gradient flow on the space of admissible coarse-grainings.

\textbf{Local update channel.} For each cell $i$, define the local leakage gradient:
\begin{equation}
\delta_i \mathcal{L}_{\text{leak}} := \left. \frac{\partial \mathcal{L}_{\text{leak}}}{\partial E_i} \right|_{E_{j \neq i} \text{ fixed}}
\end{equation}
To strictly avoid circularity with the emergent geometry $d_G$, this gradient depends only on cells $j$ within the \textbf{Algebraic Lieb-Robinson neighborhood} $N_{\text{LR}}^{\text{alg}}(i)$ of cell $i$. This neighborhood is defined intrinsically on the pre-geometric substrate using operator commutators under the canonical flow, completely independent of the correlation graph:
\begin{equation}
N_{\text{LR}}^{\text{alg}}(i) := \{ j \in V : \| [\sigma_{\Delta s}(\mathcal{A}_i), \mathcal{A}_j] \| > \epsilon \}
\end{equation}

\textbf{Causal flow.} The coarse-graining maps evolve under a local descent:
\begin{equation}
\partial_s E_i = -\eta_{\text{flow}} \cdot \Pi_{\text{adm}} \left[ \delta_i \mathcal{L}_{\text{leak}} + \lambda_{\text{drift}} \cdot \delta_i \mathcal{L}_{\text{drift}} \right]
\end{equation}

\textbf{Causality guarantee.} Because $\delta_i \mathcal{L}_{\text{leak}}$ depends only on $E_j$ for $j \in N_{\text{LR}}^{\text{alg}}(i)$, updates at cells $i$ and $k$ with non-overlapping algebraic neighborhoods \textbf{strictly commute}:
\begin{equation}
[ \partial_s E_i , \partial_s E_k ] = 0 \quad \text{whenever } N_{\text{LR}}^{\text{alg}}(i) \cap N_{\text{LR}}^{\text{alg}}(k) = \emptyset.
\end{equation}

\textbf{Definition of the selected cell net.} The projection-level cell net is the fixed point of this flow:
\begin{equation}
E^* := \lim_{s \to \infty} E(s), \quad \text{where } \left. \partial_s E_i \right|_{E=E^*} = 0 \quad \text{for all } i.
\end{equation}

\textbf{Variational interpretation status.} A local causal relaxation is a proposed mechanism. The current exact code exhaustively ranks finite partitions with sampled objectives; it does not implement the flow or prove that its attractors equal the global ranking.

Effective openness is one possible stabilization hypothesis: UV degrees of freedom
traced out by the Type I funnels may act as an environment for coarse cells. The
current CA analogy configures a cooling rule but has no matched no-cooling control,
so it establishes neither a cooling effect nor a universal thermodynamic requirement,
a Type III mechanism, or spatial collapse. Those stronger implications remain open
questions rather than consequences of the present calculation.

\subsubsection{\texorpdfstring{$\Pi_{\text{loc}}$}{Pi\_loc}: locality from correlations (graph metric from mutual information)}
Define the mutual information between cells $i$ and $j$ via \textbf{Araki relative entropy}:
\begin{equation}
I_{ij} = S_{\text{Araki}}(\omega_{i \cup j} \| \omega_i \otimes \omega_j)
\end{equation}

\textbf{Distance kernel.} Define an emergent pairwise distance by a fixed monotone decreasing map $f$:
\begin{equation}
d_{ij} = \ell_* \cdot f(I_{ij} / I_0),
\end{equation}

\textbf{Markovian Geometric Filtering \& Weighted Graph.} 
Raw mutual information can falsely identify highly entangled but geometrically distant nodes (e.g., Bell pairs, error-correcting codes) as ``adjacent.'' To ensure strict geometric proximity, the edge weight kernel applies a \textbf{Quantum Conditional Mutual Information (QCMI)} filter. If the correlations between $i$ and $j$ are entirely mediated by an intermediate neighborhood $k$, their QCMI vanishes: $I(i:j | k) \approx 0$. By restricting continuous edges to pairs with strictly non-Markovian (direct) QCMI, the framework systematically separates true geometric proximity from long-range topological entanglement. The final weighted graph is defined as $w_{ij} := \kappa(I_{ij})$ applied only to unscreened pairs.

\textbf{Graph geodesic distance.} With edge lengths $d_{ij}$, define graph-geodesic distance:
\begin{equation}
d_G(i,j) = \inf_{\text{paths } i\to j} \sum_{(a,b) \in \text{path}} d_{ab}.
\end{equation}

\subsubsection{\texorpdfstring{$\Pi_{\text{geom}}$}{Pi\_geom}: emergent geometry via complexity-stress minimization}
\label{subsec:pi-geom}
\textbf{Ontic vs. Epistemic distinction.} The physical (ontic) reality in this framework is solely the \textbf{discrete correlation graph} $(V, d_G)$. The embedding into $\mathbb{R}^{D^*}$ described below is strictly an \textbf{epistemic coordinate chart}: a lossy compression constructed by macroscopic observers to extract the effective dimension $D^*$. 

\textbf{1. Dimension-selection diagnostic.} The finite pipeline defines a selected embedding dimension $D^*$ as the integer that minimizes a declared \textbf{Complexity-Stress Functional} $F(D)$:
\begin{equation}
D^* := \underset{D \in \mathbb{N}}{\text{argmin}} \left[ \text{Stress}(D) + \lambda_{\text{dim}} |D - D_S|^2 \right]
\end{equation}
where $\text{Stress}(D)$ is the standard MDS stress, $D_S$ is the spectral-dimension diagnostic of the graph, and $\lambda_{\text{dim}}$ is an explicit tunable penalty. It can change the selected dimension in non-geometric controls and must be reported and swept; the code does not implement a $\lambda_{\text{dim}}\to0^+$ limiting procedure.

\textbf{2. Relational Embedding (Epistemic Coordinates).} Given $D^*$, the coordinates $\{x_i\}$ are determined by the stress-minimizing configuration in $\mathbb{R}^{D^*}$. 

\textbf{Local embedding status.} MDS seeks a low-stress Euclidean representation. Local flattening is an embedding property, not a derivation or test of the Einstein equivalence principle.

\textbf{Curvature status.} The current Python path uses a two-dimensional embedding-Delaunay angle-deficit proxy. Intrinsic complexes, dual volumes, Regge tensors, and refinement remain open.

\textbf{3. Local Metric Reconstruction Diagnostic.}
\begin{equation}
\min_{h \succ 0} \sum_{j \in N(i)} \left( d_G(i,j)^2 - \Delta x^T h \Delta x \right)^2  +  \lambda_{\text{spd}} \| h - I \|_F^2
\end{equation}
This is a regularized fit in MDS coordinates with explicit rank, condition, and residual diagnostics; it is not yet an intrinsic continuum metric.

\textbf{4. Metric-reconstruction breakdown diagnostic.} A large design-matrix condition number means the local fit is poorly identifiable. It is neither an event horizon nor a spacetime singularity; causal/geodesic completeness requires a separate dynamical test.

\subsubsection{\texorpdfstring{$\Pi_{\text{time}}$}{Pi\_time}: clock-rate constraint candidate}
\textbf{1. The Clock-Rate Constraint Ansatz.}
\begin{equation}
(\Delta_w + \mu^2 I) \Phi = \delta\rho
\end{equation}
where:
\begin{itemize}
    \item $\Delta_w$ is the graph Laplacian on the mutual-information weighted graph: $(\Delta_w \Phi)_i := \sum_j w_{ij}(\Phi_i - \Phi_j)$.
    \item $\delta\rho$ is the source term (defined below).
    \item $\mu$ is an optional screening parameter. On a finite connected graph, $\mu=0$ requires an explicit zero-mode compatibility policy; $\mu>0$ makes the operator invertible. Neither choice establishes a continuum $1/r$ law without refinement and boundary tests.
\end{itemize}
\textit{Status:} The solver is a tested elliptic graph constraint. Its relationship to the ADM Hamiltonian constraint remains conjectural.

\textbf{2. Source candidates and a retained falsification result.}
The original proposal was
\begin{equation}
\delta\rho_i := -\frac{1}{s_0} \left[ \bar{S}_{\text{Araki}}(\omega_{i} \| \omega_{i}^{\text{vac}}) \right]
\end{equation}
but exact tests find $D(\rho(\epsilon)\|\sigma)=O(\epsilon^2)$ and entropy dependence at fixed energy. Raw relative entropy is therefore rejected as a standalone linear source in the tested weak-response regime. For a faithful finite KMS state, $K_\sigma=\beta H+\log Z$ and $\Delta\langle K_\sigma\rangle=\beta\Delta\langle H\rangle$. The exact backend tests the separate microscopic candidate
\begin{equation}
\delta\rho_i^{(E)}=-g\,\beta\,\Delta\langle h_i\rangle,
\qquad \sum_i h_i=H.
\end{equation}
The source density sums to minus the global modular-energy change. Under the affine interpolation $\rho(\epsilon)=(1-\epsilon)\sigma+\epsilon\rho_{\mathrm{exc}}$, every fixed-observable expectation and the linear graph solve are exactly proportional to $\epsilon$; those legs are analytic-identity regressions, not falsification tests. A genuinely non-affine family $\rho(\epsilon)\propto\exp[-\beta(H+\epsilon V)]$, with localized $V=-h_{\mathrm{center}}$, has a nonzero first-order modular/energy susceptibility in the declared finite sweep ($\beta\leq3$), while nested-window fits of $D/\epsilon^2$ converge to a positive coefficient. Conversely, an isospectral local unitary family has $\Delta S=0$ and $D(\rho\|\sigma)=\Delta\langle K_\sigma\rangle=\beta\Delta\langle H\rangle=O(\epsilon^2)$. The response order is therefore family-specific. A separate finite-chain quench conserves and spreads the local-energy profile only under noncommuting dynamics. Naive one-site reduced modular energy is blind in the symmetric KMS control. The split into $h_i$ is a supplied microscopic convention, not a unique covariant stress tensor or a derived gravitational law.

\textbf{3. The Baseline Vacuum.} In the finite model, $\omega^{\text{vac}}=e^{-\beta H}/Z$ is the Gibbs/KMS reference for the chosen Hamiltonian. A representation-independent and unique physical-vacuum construction remains open.

\textbf{4. Emergent Proper Time.} The local proper time $\tau$ along a phase trajectory is then constructed by integrating the \textbf{Laplacian-derived clock rate}:
\begin{equation}
d\tau(i) = \beta_0 \cdot \exp(\Phi_i) \, dS_{\text{act}}
\end{equation}
where $dS_{\text{act}}$ is a proposed action-distance increment. The simulator currently substitutes its numerical step $dt$; an independent clock observable remains open.

\textbf{Minimal emergent spacetime line element (preferred slicing).}
\begin{equation}
ds^2 = -c^2 d\tau^2 + h_{ab}(x) dx^a dx^b,
\end{equation}

\subsection{Symbol glossary}
\begin{table}[htbp]
\caption{Symbol glossary}
\begin{ruledtabular}
\begin{tabular}{p{0.15\linewidth} p{0.8\linewidth}}
\textbf{Symbol} & \textbf{Meaning} \\
\colrule
$\mathcal{S}$ & Substrate object \\
$\mathcal{A}$ & Substrate algebra of relational observables \\
$\mathcal{A}_{\text{rel}}$ & Relational ``bulk'' algebra (pre-geometric) \\
$\mathcal{A}_F$ & Finite internal algebra encoding gauge/representation structure \\
$\omega$ & State on $\mathcal{A}$ \\
$\alpha$ & Internal symmetry action on $\mathcal{A}$ \\
$H$ & Phase/action ordering generator (not substrate time) \\
$\Pi$ & Unique physical projection map \\
$M^3$ & Emergent 3D spatial structure \\
$g_{\mu\nu}$ & Emergent spacetime metric in projection \\
$t$ & Emergent time coordinate/order metric in projection \\
$\rho(x)$ & Encoding density field in projection \\
$\ell_*$ & Minimal encoding resolution scale (Planck-like) \\
$\eta$ & Dimensionless coefficient in the area-law capacity bound \\
$\gamma$ & Junction-access gating parameter (open module) \\
$\Delta$ & Junction-access correction functional (open module) \\
\end{tabular}
\end{ruledtabular}
\end{table}

\subsection{Substrate algebra and parameter budget}
\textbf{Substrate data.} The substrate is specified by: $\mathcal{S} = (\mathcal{A}, \omega, \alpha)$. Continuous symmetry and continuous phase/action order are treated as idealized descriptions. Physical finiteness is enforced at the projection level via the area-law bound and bounded intensities.

\subsection{Parameter budget (what must be fixed or fitted)}
\begin{description}
    \item[\textbf{Structural (discrete) choices}] (Architecture decisions) Choice of algebra class for $\mathcal{A}_{\text{rel}}$; choice of internal algebra $\mathcal{A}_F$; compact internal symmetry group; equivalence-class uniqueness of projection postulate. \textit{Must be defended by conceptual minimality and consistency.}
    \item[\textbf{Empirical scale-setting constants}] (Expected to be measured) $\ell_*$ (minimal encoding patch); $\eta$ (area-law coefficient); effective constants matching $(c, \hbar, G, \Lambda)$. 
    \item[\textbf{Projection response functions}] (Primary risk of ``parameter fitting'') Mapping from encoding density to metric scaling; mapping from phase-order to clock time. \textit{Must be restricted to minimal families with clear constraints and matching requirements.}
    \item[\textbf{Substrate state selection}] Choice/characterization of $\omega$. \textit{Must be constrained by a selection principle (symmetry, equilibrium relative to canonical flow, minimal information, etc.) to avoid embedding arbitrary structure by hand.}
    \item[\textbf{Open junction module}] $\gamma(\text{regime})$ and $\Delta(\ldots)$. \textit{Must satisfy $\gamma \approx 0$ in all tested regimes; left open otherwise, but constrained by normalization/positivity and by compatibility with known no-signaling bounds.}
\end{description}

% ----------------------------------------------------------------
\section{Projection outputs and emergent time}
% ----------------------------------------------------------------
Time does not exist at the substrate level. Instead, projection constructs an objective time-order by mapping phase/action order to a temporal metric: $t = f(\Phi)$ with $f' > 0$.

% ----------------------------------------------------------------
\section{Finite distinguishability and Planck-scale resolution}
% ----------------------------------------------------------------
\textbf{1. Algebraic Cut-Capacity.} Let $R \subset V$ be a finite set of emergent funnels. Define the \textbf{cut-capacity} $\mathrm{Cap}(\partial R)$ purely from the correlation graph weights:
\begin{equation}
\mathrm{Cap}(\partial R) := \sum_{i\in R, j\notin R} \kappa(I_{ij})
\end{equation}

\textbf{2. The Araki Entropy Bound.} The distinguishable information content of $R$ is limited by the cut-capacity:
\begin{equation}
S_{\text{Araki}}( \omega|_R \| \omega_{\text{vac}}|_R ) \leq \eta \cdot \mathrm{Cap}(\partial R)
\end{equation}

\textbf{3. Area-scaling conjecture.} If the correlation graph admits a stable low-distortion three-dimensional limit, the proposed matching target is $\mathrm{Cap}(\partial R) \propto A_{\text{geo}}(\partial R) / \ell_*^2$. This is unproved and has no refinement result in the current code.

The operational consequence of the capacity bound is a limit on the effective dimension of the Type I funnel factor $N_R$:
\begin{equation}
\dim(N_R) \lesssim \exp\big( \eta \cdot \mathrm{Cap}(\partial R) \big)
\end{equation}

% ----------------------------------------------------------------
\section{Quantum statistics as projection-limited inference}
% ----------------------------------------------------------------
Measurement outcomes are computed by the Born-rule form in empirically accessed regimes: $P(i) = \text{Tr}(\rho_R E_i)$.
% ----------------------------------------------------------------
\section{Emergent geometry and GR matching}
% ----------------------------------------------------------------

\subsection{Metric definition as projection output}
The projection is intended to construct effective geometry from correlations; the current implementation contains finite embedding and local-fit diagnostics only.

\textbf{Entanglement-shear proposal.} A future model may couple clock-rate contrasts to changing correlations and spatial diagnostics. No such evolution law, shift vector, extrinsic curvature, or tensor constraint is implemented; this is not a derivation of nonlinear curvature.

\subsection{Proposed GR matching protocol (unimplemented)}
The code currently provides only a tensor-agnostic mismatch function and a two-dimensional embedding-space angle-deficit proxy.

\subsubsection{Candidate intrinsic discrete geometry}
\begin{enumerate}
\item \textbf{Intrinsic Discrete Geometry:} Construct and test a Vietoris--Rips or other justified complex from $(V,d_G)$. The current MDS-Delaunay construction is a declared proxy, not a guaranteed homeomorphic or evidence-equivalent intrinsic complex.

\item \textbf{Discrete Metric:} Assign edge squared-lengths $l_{ij}^2 = h_{ab}(x_i) \Delta x^a \Delta x^b$ consistent with the reconstructed local metric.

\item \textbf{Discrete Curvature (Regge):} Compute the curvature using \textbf{Regge Calculus}. The curvature is concentrated on the $(D^*-2)$-dimensional bones (hinges). For $D^*=3$, these are edges. The deficit angle $\epsilon_h$ at hinge $h$ is:
\begin{equation}
\epsilon_h = 2\pi - \sum_{\text{cell} \supset h} \theta_{\text{cell}}(h)
\end{equation}
where $\theta_{\text{cell}}(h)$ is the dihedral angle of the tetrahedron at hinge $h$.
\end{enumerate}

\subsubsection{Proposed closure mismatch}
If validated geometric and source tensors become available, compare them using:
\begin{equation}
G_h \cdot l_h := \epsilon_h - \Lambda V_h
\end{equation}

Define the \textbf{closure mismatch} $M(L)$ over a region $V_L$:
\begin{equation}
M(L) := \left\| G_{\text{Regge}} - 8\pi G \, \Pi_{\text{graph}}(T_{\mu\nu}^{\text{eff}}) \right\|_L
\end{equation}
Neither $G_{\text{Regge}}$, a conserved $T_{\mu\nu}^{\text{eff}}$, nor $\Pi_{\text{graph}}$ is implemented. Raw Araki contrast fails the tested linear-source criterion, and scalar clock data cannot replace the missing tensor components.

\textbf{Weak-field consistency target.} A future calibrated model must compare its independently specified source with the discrete potential under refinement:
\begin{equation}
4\pi G \rho_i \approx (\Delta_{\text{VR}} \Phi)_i
\end{equation}

% ----------------------------------------------------------------
\section{Standard Model compatibility (program)}
% ----------------------------------------------------------------
A commonly studied finite algebra whose unitary structure can realize the Standard Model gauge group is: $\mathcal{A}_F = \mathbb{C} \oplus \mathbb{H} \oplus M_3(\mathbb{C})$. Treat each effective cell as carrying an internal factor: $\mathcal{A}_i \approx M_{d_i}(\mathbb{C}) \otimes \mathcal{A}_F$.

% ----------------------------------------------------------------
\section{Entanglement and junction accessibility (open module)}
% ----------------------------------------------------------------
For accessible operations and regimes: $P(b | x, y) = P(b | x)$ (within experimental bounds).
To keep the possibility of additional operational access open without asserting it, introduce a gating parameter $\gamma(\text{regime}) \in [0, 1]$ and write:
\begin{equation}
P(b | x, y) = \text{Tr}(\rho_B E_b^{(x)}) + \gamma \cdot \Delta_b(x, y; \rho_{AB})
\end{equation}

% ----------------------------------------------------------------
\section{Evaluation, falsification, and test program}
% ----------------------------------------------------------------

\subsection{Mathematical status (definition / theorem / conjecture / matching)}

\begin{description}
    \item[\textbf{Substrate triple} $(\mathcal{A},\omega,\alpha)$] \textit{(Definition)} The substrate is specified as an atemporal, pre-geometric relational algebra with a compact internal symmetry action. \textbf{Failure mode:} Pure definition; failure only if later steps implicitly add substrate spacetime structure.
    \item[\textbf{Phase/action ordering flow} $\sigma_s$] \textit{(Definition / Conjecture)} A distinguished ordering flow exists; $s$ is an order parameter, not time. \textbf{Failure mode:} If derived (canonical/modular flow): show existence conditions on $(\mathcal{A},\omega)$. If postulated: keep minimal and symmetry-compatible.
    \item[\textbf{Cell net} $\{\mathcal{A}_i,E_i\}$] \textit{(Definition)} Projection-resolution stage outputs effective ``cells'' (Type I approximants / toy matrix factors) and coarse-graining maps. \textbf{Failure mode:} Toy models: explicit construction. General AQFT: implement via split-property funnels; failure if no consistent finite-resolution approximation exists.
    \item[\textbf{Symmetry-commutation}] \textit{(Constraint)} $E_i\circ\alpha_g=\alpha_g\circ E_i$ for all $g\in \text{SU}(2)$. \textbf{Failure mode:} Check directly in constructions; failure if recovered physics requires symmetry-breaking at the projection-selection level.
    \item[\textbf{Stability selection} $\mathcal{L}_{\text{leak}}$] \textit{(Definition)} Among admissible coarse-grainings, select those minimizing phase-flow leakage across cell boundaries. \textbf{Failure mode:} Toy models: compute minimizers numerically; failure if the minimizer is always trivial (information-destroying) despite retention constraints.
    \item[\textbf{Existence of minimizers}] \textit{(Theorem / Conjecture)} Minimizers exist for finite toy systems and well-posed admissible sets; general operator-algebraic existence is open. \textbf{Failure mode:} Toy models: compactness/finite search; general case: requires technical assumptions on admissible channel set and topology.
    \item[\textbf{Capacity bound} $\mathrm{Cap}(\partial R)$] \textit{(Definition)} Finite distinguishability is bounded by a non-geometric boundary-capacity functional computed from correlation weights. \textbf{Failure mode:} Pure definition; failure only if it cannot be made compatible with recovered continuum behavior.
    \item[\textbf{Emergent ``area law'' scaling}] \textit{(Matching output)} In manifold-like regimes, $\mathrm{Cap}(\partial R)\propto A_{\text{geo}}(\partial R)/\ell_*^2$. \textbf{Failure mode:} Validate in toy models that yield low-stress embeddings; falsified if scaling is generically volume-like even in 3D-like regimes.
    \item[\textbf{Emergent locality} $I_{ij}, d_G$] \textit{(Definition)} Locality is defined from QCMI-screened correlations and multi-hop routing on weighted graphs. \textbf{Failure mode:} Definition is checkable; failure mode is pathological graphs (disconnectedness, non-manifold structure) for physically relevant states.
    \item[\textbf{Dimension} $D^*$] \textit{(Finite benchmark + Matching)} Selected chain/grid cases yield $D^*=1,2$ with encoded interaction graphs. Robustness, controls, and refinement remain open.
    \item[\textbf{Spatial metric} $h_{ab}(x)$] \textit{(Embedding-coordinate prototype)} Regularized SPD fits report rank, condition, and residual but are not intrinsic continuum metrics.
    \item[\textbf{Clock-rate field} $\beta$ \textbf{and time} $\tau$] \textit{(Definition / diagnostic)} The graph solve and sign convention are implemented; a physical source and independent clock are open.
    \item[\textbf{Spacetime metric} $g_{\mu\nu}$] \textit{(Proposed construction)} Shift, extrinsic curvature, tensor constraints, and evolution are unimplemented.
    \item[\textbf{GR closure mismatch} $M(L)$] \textit{(Definition + Matching)} The discrete Regge Einstein tensor $G_{\text{Regge}}$ on the simplicial complex must satisfy Einstein closure within tolerance. \textbf{Failure mode:} Falsified if no small-parameter instantiation yields small $M(L)$ across standard tests.
    \item[\textbf{Born-rule statistics}] \textit{(Matching)} Outcome statistics must match standard QM in all currently tested regimes. \textbf{Failure mode:} Falsified by any predicted deviation already excluded; derivation from projection-limited inference is an open program item.
    \item[\textbf{Operational no-signaling}] \textit{(Matching)} Local marginals must not depend on distant choices (junction gate $\gamma\approx 0$). \textbf{Failure mode:} Falsified by any proposal that enables controllable signaling in ordinary Bell-test regimes.
    \item[\textbf{Junction-access deviations}] \textit{(Open module)} Whether additional operational access exists in extreme regimes is left open and parameterized. \textbf{Failure mode:} Must remain consistent with existing constraints; becomes testable if a concrete activation regime is specified.
    \item[\textbf{Local Lorentz covariance}] \textit{(Matching / conjecture)} SU(2) tests spatial rotations only. Dispersion, common-speed, anisotropy, and boost-sensitive tests are required.
\end{description}

\subsection{Internal validation tests (toy-model / simulation tests)}
These tests evaluate finite pieces of the proposed projection. T1 and T2 are selected toy benchmarks, T4 is a graph-solver/sign diagnostic with a manual source, and T5 is a phenomenological cellular-automaton analogy. None constitutes completion of the full projection. T3 remains open.

\textbf{T1 --- Recovery of encoded locality in selected states. [FINITE BENCHMARKS]}
Two exact simulations benchmark recovery of interaction topology encoded in their Hamiltonians and finite embedding-dimension selection. They do not validate the full projection or source physics.
\begin{itemize}
    \item \textbf{T1a --- 1D Chain.} The selected benchmark recovers encoded chain ordering and selects $D^*=1$. With four cells, the inferred three-edge chain is exactly the minimum spanning tree, so its perfect edge score is non-discriminating.
    \item \textbf{T1b --- 2D Grid.} The blind inference recovers the 12 encoded nearest-neighbor grid edges and selects $D^*=2$ for the declared parameters. Its SU(2)-invariant one-site entropy placeholder is constant and yields $\Phi=0$ after zero-mode removal; this is a recorded degeneracy, not a clock validation.
\end{itemize}

\textbf{T2 --- Stability under phase-flow. [FINITE TOY BENCHMARK]}
An exhaustive search evaluated all 105 four-cell partitions and selected contiguous 1D blocks in the tested configuration. The reported leakage is an unnormalized common-Haar-probe mean proportional to the Hilbert--Schmidt channel norm only at fixed dimension; the omitted $d(d+1)$ factor precludes cross-dimension comparison. The nearest-neighbor Hamiltonian already encodes the chain interaction graph, and the code ranks partitions rather than implementing the proposed causal gradient flow. This supports recovery of encoded locality in one small system, not topology-free emergence.

\textbf{T4 --- Graph-potential and redshift-sign diagnostic. [NUMERICAL CHECK; PHYSICAL SOURCE OPEN]}
A manually chosen negative point source was injected into a 2D MI-weighted graph. The solve preserves lattice symmetry, orders the available shells monotonically, and obeys the chosen redshift-sign convention $1+z=\exp(\Phi_{\mathrm{observer}}-\Phi_{\mathrm{emitter}})$. The graph is too small to establish a $1/r$ law, and this check does not validate a physical source or Newtonian limit.

\textbf{T5 --- Phenomenological Stability and Radiative Cooling. [CA ANALOGY]}
Only one cooling-enabled run is committed. It retains a nonzero population but declines from 33 to 27 cells and fails its stable-or-growing population criterion. No matched no-cooling artifact or multi-seed sweep exists, so the run does not establish that cooling is effective or necessary and does not validate emergent geometry.

\subsection{Distinctive empirical commitments and lethal falsifiers}
\textbf{F1 --- Dimensional Collapse.} The framework fails if controlled families systematically select a trivial or divergent dimension instead of a stable continuum value. Current finite chain/grid benchmarks are too small to establish stability or a continuum result. \\
\textbf{F2 --- Volume-Law Saturation.} Falsified if split-property funnel simulations show cut-capacity scaling with volume rather than boundary cuts in low-distortion regimes. \\
\textbf{F3 --- Lorentz-Recovery Failure.} The conjecture fails if a specified large-system limit retains direction-dependent dispersion, unequal limiting speeds, or preferred-frame observables incompatible with current bounds. The project has not derived a suppression law or numerical experimental prediction. \\
\textbf{F4 --- The Two-Potential Disconnect.} Falsified if matching Shapiro delay and spatial lensing for the same source requires fundamentally different parameter calibrations. \\
\textbf{F5 --- Junction Accessibility.} Any asserted nonzero accessible junction gate must remain compatible with no-signaling constraints; a robust marginal dependence in ordinary regimes would exclude that variant.

\begin{table*}[htbp]
\caption{A practical test matrix}
\begin{ruledtabular}
\begin{tabular}{llp{4.5cm}p{4.5cm}}
\textbf{Target claim/module} & \textbf{Observable proxy} & \textbf{What must hold} & \textbf{What would disprove it} \\
\colrule
Stability-selected locality & Leakage/drift metrics & $\mathcal{L}_{\text{leak}}$ small, drift small at fixed capacity & No stable decomposition exists without collapsing capacity \\
Emergent 3D space & Dimensionality & Best-fit embedding dimension $\approx 3$ across scales & No stable low-dimensional embedding; dimension drifts wildly \\
GR closure in weak field & Extracted $\Phi$, redshift, lensing & Matches standard weak-field phenomenology & Systematic mismatch not removable by universal calibration \\
Finite distinguishability & Scaling of $S(R)$ with boundary cut & $S(R)$ bounded by area scaling (operationally) & Demonstrated generic violation of area scaling at fixed regime \\
Junction module ($\gamma$) & Marginal dependence tests & No measurable marginal dependence in known regimes & Robust marginal dependence without classical side-channel explanation \\
\end{tabular}
\end{ruledtabular}
\end{table*}

% ----------------------------------------------------------------
\section{Discussion and open problems}
% ----------------------------------------------------------------
Open problems include:
\begin{enumerate}
    \item \textbf{Toy-model realizations of $\Pi$.} Scaling the tensor networks and GPUs to $D^* \approx 3$.
    \item \textbf{GR matching as an emergent closure.}
    \item \textbf{Dynamic cell net / dynamic spacetime (``geometrogenesis'').}
    \item \textbf{Speculative application modules.} Physical-AI, neuromorphic, and junction-based computing interpretations are outside the core gravity argument. No zero-latency communication or autonomous learning result is implemented; future studies must be separate and obey ordinary no-signaling constraints.
\end{enumerate}

% ----------------------------------------------------------------
\section{Proposed tests for extracting standard observables}
% ----------------------------------------------------------------

Only the finite graph clock-sign diagnostic is currently implemented. The transport,
intrinsic metric, geodesic, lensing, delay, and cosmology constructions below remain
proposals.

\subsection{Slicing, gauge, and a proposed Shift Vector via Optimal Transport}
\subsubsection{The Shift Vector from Intrinsic Optimal Transport}
\begin{equation}
v_{ij}^a \approx \text{Log}_{x_i}(x_j)
\end{equation}
\begin{equation}
N^a(x_i) \Delta s := \sum_j \frac{\gamma_{ij}^*}{\mu_i} v_{ij}^a
\end{equation}

\subsubsection{Discrete Weak-field diagnostic potentials}
\textbf{Time potential} $\Phi_t$ at node $i$: $\Phi_t(i) := \Phi_i - \Phi_B$. 
The discrete gravitational redshift:
\begin{equation}
1 + z = \exp(\Phi_B - \Phi_A)
\end{equation}

\textbf{Space potential} $\Phi_s$ from discrete conformal factors:
\begin{equation}
a_s(i) := \left( \frac{\det(h_{ab}(i))}{\det(h_B)} \right)^{1/6}
\end{equation}

\subsection{Proposed test A --- isolated spherical mass (weak-field, static)}
\subsubsection{Discrete Geodesic Shooting}
\begin{enumerate}
    \item \textbf{Null Geodesics:}
    \begin{equation}
    L_{\text{opt}} = \sum_{\text{edge}=(i,j)} \frac{\ell_{ij}}{\beta_{\text{avg}}(i,j)}
    \end{equation}
    \item \textbf{Redshift:}
    \begin{equation}
    1+z = \frac{\nu_{\mathrm{emit}}}{\nu_{\mathrm{obs}}}
    = \frac{\beta(B)}{\beta(A)} = \frac{\exp(\Phi_B)}{\exp(\Phi_A)}
    \end{equation}
    \item \textbf{Shapiro Delay:}
    \begin{equation}
    \Delta \tau \approx \sum_{\text{edges}} \ell_{ij} \big(1 - \exp(\Phi_{\text{avg}})\big)
    \end{equation}
\end{enumerate}

\subsection{Proposed test C --- homogeneous cosmology patch (FLRW-like)}
\subsubsection{Discrete Scale Factor}
\begin{equation}
V_P(t_k) := \sum_{T \in \mathcal{T}_{\text{VR}}} \text{Vol}(T)
\end{equation}
\begin{equation}
a(t_k) := \left( \frac{V_P(t_k)}{V_P(t_0)} \right)^{1/3}
\end{equation}

\subsubsection{Hubble Rate}
\begin{equation}
H(t_k) \approx \frac{1}{a(t_k)} \frac{a(t_{k+1}) - a(t_k)}{\Delta \bar{\tau}_k}
\end{equation}

\subsubsection{Failure Modes (Falsification)}
\begin{itemize}
    \item \textbf{Volume Collapse:} If the Vietoris-Rips complex degenerates (slivers, zero-volume tetrahedra) despite a growing $V_P$, the geometry is not manifold-like.
    \item \textbf{Anisotropy:} If $a(t)$ differs significantly when measured along different graph axes (using discrete direction-dependent correlators), the emergent space is not FLRW.
\end{itemize}

\subsection{How these proposed tests connect to falsification}
These examples define \textit{what to compute} from a candidate instantiation $(\mathcal{A}, \omega)$ and the specified selection rule $E$. The framework is disfavored if, under its own constraints (symmetry commutation, retention, area capacity, and stability), it cannot realize:
\begin{itemize}
    \item a weak-field isolated-mass regime reproducing redshift and lensing;
    \item multi-lens weak-field behavior with approximate superposition and shear;
    \item a homogeneous expanding patch with robust scale-factor diagnostics.
\end{itemize}
In other words: these examples translate ``emergent geometry'' from a slogan into a set of explicit computational tests.

\end{document}
